\clearpage
\section{Phase Noise}
\label{sec:phase-noise}
\raggedbottom% this section's pages end naturally (flushbottom spread large gaps through the step list); restored at the end
% Generated by FR3_Characterization/plot_datasheet_pn.py -- do not edit by hand.
\newcommand{\pnLoLow}{8}
\newcommand{\pnLoHigh}{22}
\newcommand{\pnScaleExpected}{8.8}
\newcommand{\pnScaleMin}{8.0}
\newcommand{\pnScaleMax}{8.9}
\newcommand{\pnRmsLow}{0.82}
\newcommand{\pnRmsHigh}{2.13}
\newcommand{\pnCount}{8}
\newcommand{\pnUnwantedMin}{20}
\newcommand{\pnUnwantedMax}{65}
\newcommand{\pnOctoFig}{images/sec14_fig9.pdf}
\newcommand{\pnOctoRms}{0.24}
\newcommand{\pnBoardTenRms}{1.02}
\newcommand{\pnOctoRmsRatio}{4.3}
\newcommand{\pnOctoGainMin}{12}
\newcommand{\pnOctoGainMax}{21}


\textbf{A Motivating Example.} An ideal LO would be a perfect sine wave: all of its power at exactly one
frequency, a single infinitely thin line on a spectrum analyzer. A real LO wobbles slightly in phase from moment
to moment, and that wobble spreads a little of its power out into ``skirts'' on either side of the carrier.
\emph{Phase noise} measures how much power sits in those skirts, at a given distance (the \emph{offset}) from
the carrier, in a 1\,Hz bandwidth, relative to the carrier itself --- hence the unit dBc/Hz. For example,
−90\,dBc/Hz at a 100\,kHz offset means that 100\,kHz away from the carrier, each 1\,Hz slice holds a
billionth of the carrier's power. Lower is better.

\textbf{Why it matters.} Every signal the unit up- or down-converts is multiplied by the LO, so the LO's
skirts are copied onto every signal that passes through. On the transmit side they blur each transmitted
subcarrier into its neighbors and smear the constellation of a modulated signal. On the receive side a strong
signal on a nearby channel, mixed with the LO's skirts, lands on top of a weak wanted signal (``reciprocal
mixing''). A single summary figure is the \emph{RMS phase error}: the total phase wobble, in degrees, integrated
over all the measured offsets. Higher-order modulations (64-QAM and above) need it to be a fraction of a
degree to a few degrees, depending on the signal.

\textbf{How it's measured.} The LO is not brought out to a connector of its own, so it is measured through the
small amount of it that leaks out of the ``TX-side FR3 RF Out'' connector, with nothing connected to the
``TX-side FR1 IF In'' connector except a 50\,$\Omega$ termination. The leakage is a copy of the LO, so its
phase noise is the LO's phase noise. Three steps make the measurement clean:\footnote{Measured with a Keysight
UXA signal analyzer and its Phase Noise application (log-plot measurement, 100\,Hz to 1\,MHz offset, signal
tracking on, 10 averages). The analyzer's own phase noise is well below the unit's at every offset shown.}
\begingroup\setlength{\leftmargini}{2.2em}\setlength{\topsep}{0.4em}% indent slightly; fixed (non-stretching) spacing
\begin{enumerate}\setlength{\itemsep}{0.4em}\setlength{\parsep}{0pt}
  \item \textbf{A fixed, moderate carrier level.} The I/Q DC bias described in Section~\ref{sec:lo-suppression}
        is moved slightly away from its LO-null point until the leakage at the ``TX-side FR3 RF Out'' connector
        is exactly 0\,dBm --- strong enough to measure well, and well below the level at which the TX chain
        starts to compress.
  \item \textbf{A filter that passes only the LO.} The LMX2820 synthesizer produces one other tone besides the
        LO. Up to 11.3\,GHz its oscillator drives the output directly, so its second harmonic (twice the LO
        frequency) is present; above 11.3\,GHz the output is the oscillator's frequency doubled, so the
        oscillator's own frequency (half the LO frequency) is present. At each LO frequency the AM3153 filter
        is set to pass the LO and block that other tone, so the analyzer measures the LO and nothing
        else.\footnote{The filter is chosen with the same two-step rule as the ``optimal'' filter of
        Section~\ref{sec:sideband-suppression}: keep the LO within 3\,dB of the best available, then block the
        unwanted tone as much as possible. With the filter in place, the unwanted tone sat between
        \pnUnwantedMin{} and \pnUnwantedMax\,dB below the LO at the ``TX-side FR3 RF Out'' connector.}
  \item \textbf{Time to settle.} After each retune, the synthesizer is given 5\,s to settle fully before the
        measurement starts.
\end{enumerate}
\endgroup

\textbf{What's plotted.} \pnCount{} LO frequencies, from \pnLoLow{} to \pnLoHigh\,GHz in 2\,GHz steps, each in
its own figure: phase noise against offset from the carrier, from 100\,Hz to 1\,MHz. Every curve is the average of
10 complete measurements. Each figure shows that average twice: the thin line is the average as measured
(unsmoothed), and the thick line is the analyzer's lightly smoothed version of it, which makes the trend easier
to read. The figures follow on the next page; the values marked at every decade offset, the table after them, and
the RMS phase error all come from the smoothed curve.\footnote{All measurements use the unit's own
10\,MHz TCXO and internal LMX2820 synthesizer, the default configuration.}

\clearpage
% Page of its own: all eight per-LO figures (user layout, 2026-10-04). Section-local, slightly narrower figpair
% so four rows fit on one page; scoped so the document-wide \figpair is unchanged.
\begingroup
\renewcommand{\figpair}[6]{%
  \noindent
  \begin{minipage}[t]{0.44\linewidth}\centering\includegraphics[width=\linewidth]{#1}\captionof{figure}{#2}\label{#3}\end{minipage}\hfill
  \begin{minipage}[t]{0.44\linewidth}\centering\includegraphics[width=\linewidth]{#4}\captionof{figure}{#5}\label{#6}\end{minipage}
  \par\vspace{0.6em}
}
\figpair{images/sec14_fig1.pdf}{LO = 8\,GHz.}{fig:14-1}%
        {images/sec14_fig2.pdf}{LO = 10\,GHz.}{fig:14-2}
\figpair{images/sec14_fig3.pdf}{LO = 12\,GHz.}{fig:14-3}%
        {images/sec14_fig4.pdf}{LO = 14\,GHz.}{fig:14-4}
\figpair{images/sec14_fig5.pdf}{LO = 16\,GHz.}{fig:14-5}%
        {images/sec14_fig6.pdf}{LO = 18\,GHz.}{fig:14-6}
\figpair{images/sec14_fig7.pdf}{LO = 20\,GHz.}{fig:14-7}%
        {images/sec14_fig8.pdf}{LO = 22\,GHz.}{fig:14-8}
\endgroup
\clearpage

\par\medskip\noindent\begin{minipage}{\linewidth}% caption and table stay on one page
\begin{center}
\captionof{table}{Phase noise of the internal LO at each decade offset, read from the smoothed curves of
Figures~\ref{fig:14-1}--\ref{fig:14-8}, and the RMS phase error integrated from 100\,Hz to 1\,MHz.}\label{tab:14-1}
\vspace{0.4em}
% Generated by FR3_Characterization/plot_datasheet_pn.py -- do not edit by hand.
\begin{tabular}{|l|ccccc|c|}
\hline
\textbf{LO} & \textbf{100\,Hz} & \textbf{1\,kHz} & \textbf{10\,kHz} & \textbf{100\,kHz} & \textbf{1\,MHz} & \textbf{RMS phase error} \\
 & \multicolumn{5}{c|}{(dBc/Hz)} & (degrees) \\
\hline
8\,GHz & −63.9 & −85.4 & −88.4 & −97.3 & −123.7 & 0.82 \\
10\,GHz & −62.3 & −83.3 & −86.5 & −95.1 & −122.1 & 1.02 \\
12\,GHz & −60.6 & −81.8 & −85.0 & −93.2 & −119.3 & 1.25 \\
14\,GHz & −58.4 & −80.9 & −83.8 & −92.0 & −118.2 & 1.46 \\
16\,GHz & −58.0 & −79.2 & −82.5 & −91.1 & −118.5 & 1.63 \\
18\,GHz & −57.4 & −78.1 & −82.0 & −90.2 & −116.9 & 1.80 \\
20\,GHz & −55.2 & −77.1 & −80.2 & −89.2 & −116.0 & 2.06 \\
22\,GHz & −55.4 & −76.5 & −80.2 & −89.3 & −117.2 & 2.13 \\
\hline
\end{tabular}

\end{center}
\end{minipage}\par

\textbf{Reading the curves.} Three features are worth noticing. First, every curve has the same shape, and
the whole curve moves up as the LO frequency goes up: from \pnLoLow{} to \pnLoHigh\,GHz the phase noise rises
by \pnScaleMin{} to \pnScaleMax\,dB at every offset. This is expected: multiplying a frequency by N multiplies its
phase wobble by N too, which raises phase noise by 20·log\textsubscript{10}(N), and
20·log\textsubscript{10}(\pnLoHigh/\pnLoLow) = \pnScaleExpected\,dB. The same rule sets the close-in part of each
curve: the LO is locked to a 10\,MHz reference, so close to the carrier it inherits that reference's phase noise
multiplied up by the ratio of the two frequencies --- at 10\,GHz, a factor of 1000, or +60\,dB. Second, beyond
roughly 200\,kHz every curve falls steeply: that is where the synthesizer's control loop stops following the
reference and the curve takes on the oscillator's own, falling, phase noise. Third, the RMS phase error grows
with frequency for the same reason, from \pnRmsLow\textdegree{} at \pnLoLow\,GHz to \pnRmsHigh\textdegree{} at
\pnLoHigh\,GHz. In practice: for the lowest phase noise, use the lowest LO frequency your FR1 and FR3 choices
allow.

\textbf{A cleaner LO: the OctoLO.} When phase noise matters most, the unit's own synthesizer can be bypassed
entirely: selecting External LO (Section~\ref{sec:hw-clock-lo}) routes a signal from the ``LO In'' connector to
both mixers instead. Figure~\ref{fig:14-9} compares the internal LO at 10\,GHz with the Pi-Radio OctoLO tuned to
the same frequency, measured with identical analyzer settings.\footnote{The OctoLO was connected directly to the
analyzer, so this figure shows the source itself. Near 1\,kHz its phase noise comes close to the analyzer's own,
so its true value there may be slightly lower still.} From 100\,Hz to 100\,kHz the OctoLO is \pnOctoGainMin{} to
\pnOctoGainMax\,dB quieter, and its RMS phase error is \pnOctoRms\textdegree{} instead of
\pnBoardTenRms\textdegree{} --- about \pnOctoRmsRatio{} times lower.

\figrow{\pnOctoFig}{LO = 10\,GHz: the internal LO compared with the Pi-Radio OctoLO. Thick lines are smoothed,
thin lines unsmoothed; both are the average of 10 measurements.}{fig:14-9}

\clearpage % start this discussion on a fresh page (it otherwise left a single line at the foot of page 38)
\textbf{What RMS phase error means for a modulation.} Picture a constellation diagram. Every symbol the
receiver gets is a dot, and the receiver decides which symbol was sent by checking which ideal point the dot is
closest to. Phase noise makes the whole constellation twist back and forth around its center, like a wheel being
jiggled: every dot slides sideways along a circle, and the RMS phase error is how far, in degrees, that twist
typically goes. A twist only causes an error once it carries a dot closer to a neighboring point than to its
own. BPSK has just two points, on opposite sides of the circle, so it shrugs off large twists; QPSK has four, a
quarter turn apart; 16-QAM, 64-QAM, and 256-QAM pack ever more points into the same space, with their outer
points crowded closest together around the circle, so each step up tolerates less twist. Wireless standards
such as 5G NR capture this by limiting how far, on average, the received dots may stray from their ideal positions --- the
error vector magnitude (EVM), a percentage of the signal's size --- with a tighter limit for each denser
modulation. Translating each limit into degrees of twist gives the most RMS phase error each modulation could
tolerate if phase noise were the \emph{only} thing disturbing it:
\par\medskip\noindent\begin{minipage}{\linewidth}% caption and table stay on one page
\begin{center}
\captionof{table}{EVM limit of each modulation, and the RMS phase error it would allow if phase noise were the only impairment.}\label{tab:evm-phase}
\vspace{0.4em}
\begin{tabular}{lcc}
\toprule
\textbf{Modulation} & \textbf{EVM limit} & \textbf{RMS phase error, if it were the only impairment} \\
\midrule
QPSK & 17.5\% & about 10\textdegree \\
16-QAM & 12.5\% & about 7\textdegree \\
64-QAM & 8\% & about 4.6\textdegree \\
256-QAM & 3.5\% & about 2.0\textdegree \\
\bottomrule
\end{tabular}
\end{center}
\end{minipage}\par
In a real link, phase noise must share that budget with everything else (thermal noise, nonlinearity, I/Q
imbalance), and both ends of the link contribute their own LO's phase noise, so a comfortable design keeps phase
noise to well under half of these limits. Against that yardstick, the internal LO (\pnRmsLow\textdegree{} at
\pnLoLow\,GHz up to \pnRmsHigh\textdegree{} at \pnLoHigh\,GHz) supports QPSK to 64-QAM comfortably at every LO
frequency, while 256-QAM is within reach only at the lower LO frequencies and becomes tight toward the top of the
range. The OctoLO, at \pnOctoRms\textdegree{}, leaves ample margin even for 256-QAM.\footnote{The RMS phase error
here is integrated from 100\,Hz to 1\,MHz offset. OFDM receivers track and remove much of the slow, close-in phase
wobble, which makes the close-in part matter less; phase noise beyond 1\,MHz, not included here, adds a little.
The comparison above is a guide to where each LO's limits lie, not a link budget.}

\textbf{An Acknowledgement.} Pi-Radio does not own phase noise measurement equipment. The measurements in this
section were made for us by Prof.\ Hamed Rahmani at New York University, on his lab's Keysight N9042B UXA signal
analyzer. We are grateful for his help. The analyzer's phase noise measurement software was a donation from
Keysight, and we are grateful to Keysight too.

\clearpage
\flushbottom
